This part of the exam consists of 7 problems located in separate exam
stations. You have 5 minutes to complete each problem. (The number of problems
was to be reduced to 5 if the night observation was successful.)

\begin{description}
\item[(P1)] \textbf{Naked eye observation from real sky with panoramic
  360-degree image.} Estimate the LST (Local Sidereal Time) at the time the
  image was taken, rounded to the nearest hour.

  Included: a panoramic 360-degree image of the sky at night at an unknown
  location; computer screen; keypad to pan around the image; coordinates of
  bright stars.
  % FIGURE NEEDED: interactive panoramic 360-degree night-sky image (on-site
  % computer station)

  \begin{center}
  \begin{tabular}{llll}
  Name & Bayer Designation & Declination (Dec) & Right Ascension (RA) \\
  \hline
  Rigil Kentaurus & $\alpha$ Cen & $-60^\circ 50' 02.3737''$
    & $14^{\mathrm{h}}39^{\mathrm{m}}36.5^{\mathrm{s}}$ \\
  Arcturus & $\alpha$ Boo & $+19^\circ 10' 56''$
    & $14^{\mathrm{h}}15^{\mathrm{m}}39.7^{\mathrm{s}}$ \\
  Vega & $\alpha$ Lyr & $+38^\circ 47' 01''$
    & $18^{\mathrm{h}}36^{\mathrm{m}}56.3^{\mathrm{s}}$ \\
  Capella & $\alpha$ Aur & $+45^\circ 59' 53''$
    & $05^{\mathrm{h}}16^{\mathrm{m}}41.4^{\mathrm{s}}$ \\
  Altair & $\alpha$ Aql & $+08^\circ 52' 06''$
    & $19^{\mathrm{h}}50^{\mathrm{m}}47.0^{\mathrm{s}}$ \\
  Aldebaran & $\alpha$ Tau & $+16^\circ 30' 33''$
    & $04^{\mathrm{h}}35^{\mathrm{m}}55.2^{\mathrm{s}}$ \\
  Antares & $\alpha$ Sco & $-26^\circ 25' 55''$
    & $16^{\mathrm{h}}29^{\mathrm{m}}24.5^{\mathrm{s}}$ \\
  Spica & $\alpha$ Vir & $-11^\circ 09' 41''$
    & $13^{\mathrm{h}}25^{\mathrm{m}}11.6^{\mathrm{s}}$ \\
  Deneb & $\alpha$ Cyg & $+45^\circ 16' 49''$
    & $20^{\mathrm{h}}41^{\mathrm{m}}25.9^{\mathrm{s}}$ \\
  Dubhe & $\alpha$ UMa & $+61^\circ 45' 04''$
    & $11^{\mathrm{h}}03^{\mathrm{m}}43.7^{\mathrm{s}}$ \\
  Polaris & $\alpha$ UMi & $+89^\circ 15' 51''$
    & $02^{\mathrm{h}}31^{\mathrm{m}}49.1^{\mathrm{s}}$ \\
  Alpheratz & $\alpha$ And & $+29^\circ 05' 26''$
    & $00^{\mathrm{h}}08^{\mathrm{m}}23.3^{\mathrm{s}}$ \\
  Schedar & $\alpha$ Cas & $+56^\circ 32' 14''$
    & $00^{\mathrm{h}}40^{\mathrm{m}}30.4^{\mathrm{s}}$ \\
  \end{tabular}
  \end{center}

\item[(P2)] \textbf{Planet observation with real sky in panoramic 360-degree
  image.} Count the number of planets visible in this image above the horizon
  and name the constellations they are in (with IAU designations, i.e.\ Ursa
  Major or UMa).

  Included: a panoramic 360-degree image of the sky at night at an unknown
  location; computer screen; keypad to pan around the image.
  % FIGURE NEEDED: interactive panoramic 360-degree night-sky image (on-site
  % computer station)

\item[(P3)] \textbf{Analemma on another planet.} Find the obliquity (axial
  tilt) of the planet.

  Included: a generated analemma (position of a star taken from the surface of
  a planet with interval separated by the mean solar day of the planet over an
  orbital period around a star) of a fictitious planet orbiting around a star.
  The result is graphed on a paper with each major grid representing
  $5^\circ$.
  % FIGURE NEEDED: analemma of the fictitious planet on a 5-degree grid

\item[(P4)] \textbf{Exposure time from a photograph.} Estimate the exposure
  time of a given ``star trails'' image.

  Included: a ``star trails'' image that was taken by a still camera capturing
  an image over a period of time; ruler.
  % FIGURE NEEDED: star-trails photograph

\item[(P5)] \textbf{Find True North from Moon shadow.} Draw an arrow pointing
  North in the data sheet.

  Included: simulated positions of the Moon shadows of a pole at certain
  intervals in the span of a day. The observer is located in the Southern
  hemisphere at latitude $27^\circ$S. The Moon's declination that night is
  $+15^\circ$. Ruler, compass (drawing tool), geometry kit.
  % FIGURE NEEDED: data sheet with simulated Moon-shadow positions of a pole

\item[(P6)] \textbf{Find latitude from equatorial mount.} Without altering the
  polar alignment, find the observer's latitude based on a previously
  polar-aligned equatorial mount.

  Included: an equatorial-mount telescope that has already been properly
  aligned to a location in the northern celestial pole; bubble level. The
  latitude dial on the mount is covered (you may not use it).

\item[(P7)] \textbf{Precision polar alignment with equatorial mount.} Perform
  a polar alignment on the equatorial mount provided.

  Included: equatorial mount with polar scope (has not been polar aligned);
  date and time (GMT, UTC+0) of the time of performing the polar alignment;
  diagram of the sky's position at the time; a light source substituted for
  Polaris to be used for proper polar alignment (already visible in the polar
  scope); longitude of observer.

  Date and Time: 30 Aug 2017 / 23:30. \quad Longitude: $10^\circ$\,E.
  % FIGURE NEEDED: diagram of the sky's position at the given time (polar
  % scope alignment aid)
\end{description}
